% ===========================================================================
%  CAPITULO 6 — Compromissos obrigatórios da V1
% ===========================================================================
\chapter{Compromissos obrigatórios da V1}
\label{cap:v1}

Os seis requisitos abaixo já foram decididos como obrigatórios para a primeira
versão, independentemente do estado atual. Para cada um: status, diferença
para o comportamento esperado, proposta de entrega e plano de aceite.
\textbf{Responsável e data de homologação não foram fixados no repositório}
e ficam marcados como \textit{[a definir]} para o time completar; este
documento não inventa prazo.

\begin{center}
\footnotesize\renewcommand{\arraystretch}{1.35}
\begin{tabularx}{\linewidth}{@{}L{14mm}L{50mm}L{36mm}Y@{}}
\toprule
\textbf{Item} & \textbf{Compromisso} & \textbf{Status} & \textbf{Maior lacuna}\\
\midrule
V1-01 & Indicadores de adoção e uso & \parcial & filtros por período,
  bimestre e turma; último acesso; baixa participação\\
V1-02 & Pré-visualização completa & \naoimpl & o recurso inteiro\\
V1-03 & Recomendações por IA & \naoimpl & o recurso inteiro\\
V1-04 & Recomendação com o tema do setup & \naoimpl & depende de V1-03; o
  tema já está gravado\\
V1-05 & Revisão docente da avaliação & \implementado & teste e2e de ajuste em
  lote\\
V1-06 & Justificativa e auditoria & \parcial & justificativa opcional;
  histórico sem tela\\
\bottomrule
\end{tabularx}
\end{center}

% ---------------------------------------------------------------------------
\section{V1-01 · Indicadores de adoção e uso}

\begin{resposta}{V1-01}{Adoção por rede, escola, turma e bimestre}{\parcial}
\begin{campos}
  \item[Existe] Relatório da diretoria por missão e escola (turmas abertas,
    turmas encerradas, estudantes com registro), recortado por regional e
    exportável; filtro de usuários ``ativos / inativos / sem primeiro
    acesso''; períodos letivos já importados do OneRoster (T03, T10).
  \item[Diferença] Faltam: filtros por \textbf{período e bimestre}, por
    \textbf{ano/série} e por \textbf{turma}; o registro de \textbf{último
    acesso} de professores e estudantes; a contagem de atividades
    configuradas, iniciadas e concluídas; a \textbf{taxa de participação}
    (estudantes que entraram sobre matriculados na turma com missão liberada)
    e a lista das turmas e escolas com participação abaixo de um limiar; a
    série no tempo. A história H07.3.2 (``números do dia'') cobre parte disso
    e está a fazer.
  \item[Proposta] (1) Gravar data e hora do último acesso por usuário a cada
    login. (2) Um painel de adoção para os três níveis de gestão, com os
    filtros acima, usando os dados que já existem: liberações, atividades
    postas no ar, progresso e conclusão por estudante. (3) Denominador
    explícito: matriculados ativos na data, de modo que transferência e
    desligamento não inflem nem zerem a taxa. (4) Indicador ``cadastrados
    $\times$ com primeiro acesso $\times$ ativos no período''.
  \item[Responsável] Produto + Engenharia \textit{[a definir]}.
  \item[Homologação] \textit{[a definir]}.
  \item[Aceite] Rede sintética com escolas de alta e baixa participação; o
    painel deve mostrar os números esperados, por filtro, e apontar as de
    baixa participação. Suíte e2e nova por papel de gestão.
\end{campos}
\end{resposta}

% ---------------------------------------------------------------------------
\section{V1-02 · Pré-visualização completa do professor}

\begin{resposta}{V1-02}{Experimentar a missão antes da turma}{\naoimpl}
\begin{campos}
  \item[Existe] Nada que permita jogar a missão: hoje ela só abre para o
    perfil estudante (T11). Existem a ficha, o guia e a revisão item a item
    dos conteúdos gerados (P01, P04).
  \item[Diferença] O recurso inteiro. Pontos que a implementação precisa
    resolver: percorrer todas as etapas e ramificações com o conteúdo
    aprovado; permitir avançar as etapas que, na aula, dependem de comando do
    professor; voltar ao preparo para corrigir um item; e não gerar
    participação, conclusão nem evidência.
  \item[Proposta] Um modo \emph{Experimentar como estudante}, aberto a partir
    da atividade, que roda a mesma missão com uma sessão marcada como
    prévia, pertencente ao professor e fora da turma. As rotas de estudante
    reconhecem a marca e não gravam progresso, conclusão, telemetria
    pedagógica nem evidência (ou gravam num espaço separado que nenhum painel
    lê). As etapas comandadas pelo professor ficam disponíveis num controle
    da própria prévia. Um botão leva de volta ao preparo da atividade.
  \item[Limitações] As 14 missões têm rotas próprias de estudante: o custo é
    por missão, e a entrega pode ser escalonada pelas missões do primeiro
    bimestre.
  \item[Responsável] Produto + Engenharia \textit{[a definir]}.
  \item[Homologação] \textit{[a definir]}.
  \item[Aceite] Por missão: o professor percorre a prévia até o fim; depois,
    painel da aula, painel de habilidades, relatórios e exportações da turma
    continuam idênticos ao estado anterior (teste automatizado que compara os
    números antes e depois).
\end{campos}
\end{resposta}

% ---------------------------------------------------------------------------
\section{V1-03 · Recomendações pedagógicas por IA}

\begin{resposta}{V1-03}{Sugestões de intervenção}{\naoimpl}
\begin{campos}
  \item[Existe] As entradas já existem: resultado por habilidade na turma
    (atingiu, em processo, não atingiu, não avaliado), os critérios por nível
    de cada habilidade, o guia da aula (objetivos, rubricas, desdobramentos)
    e a infraestrutura de IA com quota, guardrails e medição de custo. Não há
    geração de recomendação (P08).
  \item[Diferença] O recurso inteiro e o processo pedagógico em volta dele:
    autoria e versionamento dos prompts pela Pedagogia, conjunto de
    avaliação com casos aprovados por pedagogos, e regra de apresentação que
    separe dado observado de sugestão.
  \item[Proposta] Gerar, sob demanda do professor no painel de habilidades,
    uma recomendação em formato fechado: \emph{o que os dados mostram}
    (somente contagens e critérios, copiados do painel, nunca escritos pela
    IA) e \emph{o que sugerimos} (intervenções ligadas a cada habilidade em
    dificuldade, citando o código). Prompt e diretrizes ficam no servidor,
    com versão gravada junto de cada recomendação. Nenhum dado nominal de
    estudante vai ao modelo. Consome a quota do professor.
  \item[Responsável] Pedagogia (prompts e casos de avaliação) + IA
    \textit{[a definir]}.
  \item[Homologação] \textit{[a definir]}.
  \item[Aceite] Conjunto de casos de turma sintética revisado por pedagogos,
    com critérios de aprovação (pertinência, ausência de afirmação sem
    suporte nos dados, adequação à série) e taxa mínima acordada; a
    recomendação nunca afirma número que não esteja no painel.
\end{campos}
\end{resposta}

% ---------------------------------------------------------------------------
\section{V1-04 · Recomendações considerando a contextualização}

\begin{resposta}{V1-04}{Tema do setup na recomendação}{\naoimpl}
\begin{campos}
  \item[Existe] O tema, os tópicos e o componente curricular informados pelo
    professor ao criar a atividade já ficam gravados nela (P02), assim como
    a atividade de onde veio cada resultado.
  \item[Diferença] Depende de V1-03. Falta levar esses campos à geração da
    recomendação, por atividade.
  \item[Proposta] A recomendação de V1-03 recebe, além dos resultados, o
    componente, o tema e os tópicos da atividade, e propõe intervenções no
    contexto escolhido (no exemplo, desinformação aplicada a ``vacinas'' em
    Ciências), sem alterar habilidade nem critério.
  \item[Responsável] Pedagogia + IA \textit{[a definir]}.
  \item[Homologação] \textit{[a definir]}, junto de V1-03.
  \item[Aceite] Para a mesma turma sintética com dois temas diferentes, as
    recomendações mudam de contexto e mantêm as mesmas habilidades e os
    mesmos números.
\end{campos}
\end{resposta}

% ---------------------------------------------------------------------------
\section{V1-05 · Revisão docente da avaliação}

\begin{resposta}{V1-05}{O professor tem a decisão final}{\implementado}
\begin{campos}
  \item[Existe] O professor que leciona na turma ajusta o nível de qualquer
    habilidade de um ou de vários estudantes, por atividade e rodada, a
    qualquer momento, inclusive com a missão encerrada, e pode marcar ``não
    participou''. O ajuste passa a ser a origem ``professor'' e nada o
    sobrescreve depois. Também confirma, para a turma, níveis projetados.
    Gestores só leem.
  \item[Diferença] O Mapa de Produto marca a história como parcial porque
    falta o teste e2e do ajuste em lote e do histórico (H14.4.2). Ressalva:
    \emph{zerar a turma} apaga também o ajuste vigente (\adrr{ADR-155});
    o evento do ajuste continua no histórico (V1-06).
  \item[Responsável] Produto + Pedagogia \textit{[a definir]}.
  \item[Homologação] Pronto para demonstração; falta o teste citado.
  \item[Aceite] Em homologação: o professor ajusta um nível; o painel de
    habilidades, o perfil do estudante e os relatórios passam a mostrar o
    nível do professor.
  \item[Evidência] \arq{apps/\B{}api/\B{}src/\B{}evidencias/\B{}evidencias-professor.controller.ts};
    \arq{apps/\B{}hub/\B{}components/\B{}evidencias/\B{}ajustar-selecionados.tsx}.
\end{campos}
\end{resposta}

% ---------------------------------------------------------------------------
\section{V1-06 · Justificativa e auditoria obrigatórias}

\begin{resposta}{V1-06}{Trilha da revisão}{\parcial}
\begin{campos}
  \item[Existe] Cada ajuste gera um evento permanente com o valor original,
    o valor novo, o autor e a data e hora; o banco só aceita inclusão (a
    alteração é recusada e o serviço não tem permissão de apagar), e o
    histórico sobrevive a \emph{zerar a turma}. Há um motivo por lista
    (presente sem participar, problema técnico, legenda revisada, outro).
  \item[Diferença] (1) O motivo é \textbf{opcional} (a tela oferece ``sem
    motivo'') e não há texto livre de justificativa. (2) \textbf{Nenhuma tela
    ou rota consulta o histórico}: hoje ele só é lido no banco. (3) Não há
    exportação da trilha.
  \item[Proposta] Tornar o motivo obrigatório, com texto livre obrigatório
    quando ``outro''; mostrar, por estudante e habilidade, o histórico
    completo (original, cada alteração, autor, data e justificativa) ao
    professor e à coordenação; exportação da trilha para a gestão.
  \item[Responsável] Produto + Engenharia/Segurança \textit{[a definir]}.
  \item[Homologação] \textit{[a definir]}.
  \item[Aceite] Ajuste sem justificativa é recusado; após três ajustes
    sucessivos, o histórico mostra o resultado original e as três alterações
    com autor, data e motivo; tentativa de alterar ou apagar o histórico
    direto no banco com o usuário da aplicação é recusada.
\end{campos}
\end{resposta}
